\documentclass[12pt]{book}
\usepackage[papersize={6in, 9in}]{geometry}
\usepackage[utf8]{inputenc}
\usepackage[T1]{fontenc}
\usepackage[polish]{babel}
\usepackage{microtype}
\usepackage{enumitem}

\renewcommand{\baselinestretch}{1.3}

\usepackage{lettrine}
\usepackage{fix-cm}

\usepackage{fancyhdr}
\pagestyle{headings}
\fancyhf{}
\fancyfoot[LE]{\oldstylenums{\thepage}}
\fancyfoot[RO]{\oldstylenums{\thepage}}
\renewcommand{\headrulewidth}{0pt}

\newcommand{\sekcja}[1]{
  \addcontentsline{toc}{section}{#1}
  \section*{#1}
}

\newcommand{\dialog}[2]{
  \begin{itemize}[%
    ,label=#1: 
    ]
    \item #2
  \end{itemize}
}

\title{Czarne koszule i Czerwoni}
\author{Michael Parenti}

\begin{document}
\frontmatter
\maketitle

\mainmatter
\setcounter{chapter}{0}
\tableofcontents

\addcontentsline{toc}{chapter}{Wstęp}
\chapter*{Wstęp}
Ta książka jest zaproszeniem dla tych, którzy są zanurzenu w panującej ortodoksji "demokratycznego kapitalizmu" do rozważenia obrazobórczych poglądów, do zakwestionowania dogmatów mitologii wolnego rynku i wszechobecnego antykomunizmu, zarówno prawicowego, jak i lewicowego, i do rozważenia na nowo, z otwartym, lecz nie bezkrytycznym umysłem, historycznych wysiłków nieustannie oczernianych Czerwonych i innych rewolucjonistów.

Polityczna ortodoksja, odpowiedzialna ze demonizowanie komunizmu, przeszywa całe polityczne spektrum. Nawet ludzie na lewicy zinternalizowali liberalną/konserwatywną ideologię, zrównującą faszyzm i komunizm jako jednakowo złe, totalitarne bliźnięta, dwa największe masowe ruchy XX wieku. Niniejsza książka jest próbą ukazania gigantycznych różnic między faszyzmem, a komunizmem, zarówno w przeszłości, jak i współcześnie, zarówno w teorii, jak i w praktyce, szczególnie w odniesieniu do kwestii równości społecznej, akumulacji prywatnego kapitału i interesów klasowych.

Ortodoksyjna mitologia kazałaby nam również wierzyć, że zachodnie demokracje (z USA na czele) przeciwstawiały się obu totalitarnym systemom z jednakowym zapałem. W rzeczywistości przywódcy USA byli przede wszystkim oddani zapewnieniu światu bezpieczeństwa dla globalnych korporacyjnych inwestycji i systemu prywatnych zysków. W tym celu używali faszyzmu, żeby chronić kapitalizm, zapewniając, że ratują demokrację przed komunizmem.

W tej książce omawiam jak kapitalizm promuje i czerpie korzyści z faszyzmu, wartość rewolucji w postępie warunków życia, przyczyny i skutki upadku komunizmu, ciągłą aktualność marksizmu i analizy klasowej oraz bezduszną naturę władzy klasy korporacyjnej.

Ponad 100 lat temu, w swojej wielkiej pracy \emph{Les Misérables} Victor Hugo zadał pytanie: "Czy przyszłość nadejdzie?". Miał na myśli przyszłość społecznej sprawiedliwości, wolnej od "straszliwych cieni" ucisku narzucanego przez nielicznych na masach ludzkości. Ostatnio, niektórzy pisarze ogłosili "koniec historii". Twierdzą oni, że wraz z upadkiem komunizmu, monumentalna walka między alternatywnymi systemami się zakończyła. Zwycięstwo kapitalizmu jest ostateczne. Nie ma mowy o żadnych wielkich transformacjach. Globalny wolny rynek pozostanie. To co widzimy pozostanie niezmienione, teraz i na zawsze. Tym razem walka klas definitywnie się zakończyła. Zatem pytanie Hugo doczekało się odpowiedzi: przyszłość rzeczywiście nadeszła, ale niekoniecznie taka na jaką liczył.

Ta intelektualnie nijaka teoria końca historii zostałą okrzyknięta genialną egzegezą i spotkała się z hojnym przyjęciem przez kontrolowane przez korporacje media. Perfekcyjnie służy oficjalnemu spojrzeniu na świat, głosząc to, co wyższe kręgi wmawiają nam od pokoleń: że walka klasowa nie jest codzienną rzeczywistością, tylko przedawnionym pojęciem, że nieskrępowany kapitalizm pozostanie z nami już na zawsze oraz, że przyszłość należy do tych, którzy kontrolują teraźniejszość.

Ale pytanie, które rzeczywiście powinniśmy sobie zadać brzmi: czy w ogóle mamy jakąś przyszłość? Bardziej niż kiedykolwiek, gdy zagrożona jest nawet planeta, na której żyjemy, konieczne staje się zmierzenie z rzeczywistością tych, którzy grabią nasze zasoby naturalne w pogoni za nieskończonym zyskiem, tych, którzy trwonią nasze prawa i gaszą nasze wolności swojej bezkompromisowej pogoni za własnym zyskiem.

Historia uczy nas, że wszystkie elity rządzące starają się przedstawiać siebie jako naturalny i trwały porządek społeczny, nawet te, które są w poważnym kryzysie, który zagraża pochłonięciem środowiskowej bazy umożliwiającej nieustanne odtwarzanie hierarchicznej strukturę władzy i przywilejów. Wszyskie klasy rządzące lekceważą i nie tolerują odmiennych poglądów w tych sprawach.

Prawda jest niekomfortowym miejscem dla tych, którzy udają, że służą społeczeństwu, w trakcie gdy rzeczwiście służą tylko samym sobie - naszym kosztem. Mam nadzieję, że ta książka nadszarpnie to Wielkie Kłamstwo. Prawda niekoniecznie musi nas wyzwolić jak twierdzi Biblia, ale jest ważnym pierwszym krokiem w tym kierunku.

\begin{flushright}
  \textbf{--- Michael Parenti}
\end{flushright}

\chapter{Racjonalny faszyzm}
Pewnego razu, przechodząc przez Little Italy, dzielnicę Nowego Jorku, minąłem sklep z upominkami, w którym na wystawie można było zobaczyć koszulki z wizerunkiem Benito Mussoliniego wykonującego faszystowski salut. Kiedy wszedłem do sklepu i zapytałem sprzedawcy dlaczego sprzedaje takie przedmioty, odpowiedział: "Cóż, niektórym się podobają. I wiesz, być może potrzebujemy kogoś takiego jak Mussolini w tym państwie." Jego komentarz był dla mnie przypomnieniem, że faszyzm przetrwał jako coś więcej niż tylko historyczna ciekawostka.

Gorsze od plakatów lub koszulek są prace różnych pisarzy, którzy usiłowali "wyjaśnić" Hitlera, "przewartościować" Franco albo w ten czy inny sposób wybielić historię faszyzmu. We Włoszech, w latach 70. XX wieku narodził się chałupniczy przemysł książek i artykułów, zapewniających, że Mussolini nie tylko sprawił że pociągi przyjeżdżały na czas, ale także sprawił, że całe Włochy działały dobrze. Wszystkie te publikacje, wraz z wieloma opracowaniami akademickimi, mają jedną wspólną cechę: mówią niezwykle mało o ile cokolwiek o klasowej polityce faszystowskich Włoch i Nazistowskich Niemiec. Jakie podejście miały te reżimy do usług publicznych, podatków, przedsiębiorstw i warunków pracy? Na kogo korzyść działały i kosztem kogo? Większość literatury poświęconej faszyzmowi i nazizmowi nie udziela nam odpowiedzi na te pytania.\footnote{Spośród tysięcy tytułów rozprawiających o faszyzmie, jest kilka chlubnych wyjątków, które nie unikają pytań o ekonomię polityczną i władzę klasową, np. Gaetano Salvemini, \textit{Under the Ax of Fascism}; Daniel Guerin, \textit{Fascism and Big Business}; James Pool i Suzanne Pool, \textit{Who Financed Hitler}; Palmiro Togliatti, \textit{Lectures on Fascism}; Franz Neumann, \textit{Behemoth}; R. Palme Dutt, \textit{Fascism and Social Revolution}.}

\sekcja{Plutokraci wybierają autokratów}

Zacznijmy od rzucenia okiem na twórcę faszyzmu. Młodość urodzonego w 1883 roku syna kowala, Benito Mussoliniego była naznaczona ulicznymi bójkami, aresztowaniami, pobytami w więzieniu i radykalnymi działaniami politycznymi. Przed I wojną światową Mussolini był socjalistą. Jako wybitny organizator, agitator i uzdolniony dziennikarz, został redaktorem oficjalnej gazety partii socjalistycznej. Mimo tego, wielu jego towarzyszy podejrzewało go o bycie bardziej zainteresowanym własnym rozwojem niż rozwojem socjalizmu. I rzeczywiście, gdy włoska klasa wyższa kusiła go rozpoznaniem, wsparciem finansowym i obietnicami władzy, nie zawahał się zmienić strony.

Pod koniec I wojny światowej, Mussolini - socjalista, który organizował strajki dla pracowników i chłopów stał się Mussolinim - faszystą, który rozbijał strajki dla finansjerów i właścicieli ziemskich. Przy użyciu ogromnej sumy pieniędzy otrzymanych od bogatych interesantów, wypromował się na narodowej scenie jako uznany lider \textit{i fasci di combattimento}, ruchu złożonego z ubranych ubranych w czarne koszule eks-oficerów wojskowych i rozmaitych twardzieli, którym nie przyświecała żadna konkretna doktryna polityczna poza militarystycznym patriotyzmem i konserwatywną niechęcią do wszystkiego, co związane z socjalizmem i zorganizowanymi związkami zawodowymi. Faszystowskie \textit{Czarne koszule} spędzały czas atakując związkowców, socjalistów, komunistów i spółdzielni rolniczych. Po I wojnie światowej we Włoszech ukształtował się model demokracji parlamentarnej. Niskie płace rosły, a pociągi już przyjeżdżały na czas. Ale kapitalistyczna gospodarka była w powojennej recesji. Inestycje stały w miejscu, ciężki przemysł działał znacznie poniżej swoich możliwości, a korporacyjne zyski i eksport agrobiznesu malały.

Aby utrzymać swój poziom zysku, wielcy posiadacze ziemscy i industrialiści musieli obniżać płace i podnosić ceny. Państwo z kolei musiało zapewnić im olbrzymie dotacje i ulgi podatkowe. Żeby sfinansować to korporacyjne rozdawnictwo, ludność musiała zostać obłożona cięższymi podatkami, a wydatki na usługi publiczne i opiekę społeczną - znacznie obniżone. Jest to mechanizm, który dzisiaj może brzmieć dla nas znajomo.

Lecz rząd nie miał kompletnej wolności w osiąganiu tych celów. Do 1921 roku, wielu włoskich pracowników i chłopów było zrzeszonych w związkach zawodowych i organizacjach politycznych. Za pomocą demonstracji, strajków, bojkotów, przejęć fabryk i pól uprawnych wywalczyli prawo do zrzeszania się, wraz z udogodnieniami w płacach i warunkach pracy.

By narzucić z pełną surowością oszczędność na robotnikach i chłopach, rządzące interesy ekonomiczne musiałyby znieść demokratyczne prawa, które pomagały masom chronić ich już i tak skromne standardy życia. Rozwiązaniem było zniszczenie ich związków zawodowych, organizacji politycznych i wolności obywatelskich. Industrialiści i wielcy posiadacze ziemscy chcieli mieć u steru kogoś, kto mógłby złamać władzę zorganizowanych robotników i pracowników rolnych i wymusić srogi porządek na masach. Do tego zadania Mussolini, uzbrojony w swoje bandy Czarnych koszul wydawał się odpowiednim kandydatem.\footnote{Od stycznia do maja 1921 roku, "faszyści zniszczyli 120 siedzib związków zawodowych, zaatakowali 243 socjalistycznych ośrodków i innych budynków, zabili 202 robotników (oprócz 44 zabitych przez policję i żandarmerię) i okaleczyli 1144". W tym okresie czasu aresztowano 2240 robotników i tylko 162 faszystów. W latach 1921-22, aż do przejęcia państwowej władzy przez Mussoliniego, "spalono 500 hal robotniczych i sklepów spółdzielczych oraz rozwiązano 900 socjalistycznych gmin"; R. Palme Dutt, \textit{Fascism and Social Revolution}, 124.}

W 1922 roku, \textit{Federazion Industriale}, złożona z liderów przemysłu i przedstawicieli ze stowarzyszeń bankowych i agrobiznesowych, spotkała się z Mussolinim by zaplanować "Marsz na Rzym", przeznaczając na to przedsięwzięcie 20 milionów lir. Z dodatkowym wsparciem najwyszych rangą oficerów włoskiego wojska i szefów włoskiej policji, faszystowska "rewolucja" - a w rzeczywistości zamach stanu - zwyciężyła.

W ciągu 2 lat po przejęciu władzy Mussolini zamknął wszystkie opozycyjne gazety i zmiażdzył partie socjalistyczną, liberalną, demokratyczną i republikańską, które łącznie zdobyły ok. 80\% głosów. Przywódcy związkowców i chłopów, parlamentarni delegaci i inni krytycy nowego reżimu zostali pobici, zesłani na wygnanie albo zamordowani przez faszystowskie bojówki \textit{squadristi}. Włoska partia komunistyczna spotkała się z najcięższymi represjami, lecz mimo to, prowadziła odważny, podziemny opór, który z czasem przerodził się w zbrojną walkę przeciwko Czarnym koszulom i niemieckim siłom okupacyjnym.

W Niemczech pojawił się podobny schemat współpracy między faszystami, a kapitalistami. Niemieccy robotnicy i pracownicy rolni wywalczyli prawo do zrzeszania się w związkach zawodowych, ośmiogodzinny dzień pracy i ubezpieczenie od bezrobocia. Lecz żeby przywrócić dawny poziom zysków, przemysł ciężki i finansjera domagały się obniżek płac swoich dla swoich pracowników oraz ogromnych państwowych dotacji i ulg podatkowych dla siebie.

W latach 20 XX wieku nazistowskie \textit{Sturmabteilung} (SA), czyli oddziały szturmowe w brunatnych koszulach, wspierane przez biznes, były wykorzystywane głównie jako paramilitarna siła antyrobotnicza, której zadaniem było terroryzowanie robotników i pracowników rolnych. Do roku 1930, większość potentatów doszła do wniosku, że Republika Weimarska nie spełnia już ich wymagań i jest zbyt uległa dla klasy robotniczej. Znacznie zwiększyli środki przeznaczone na finansowanie Hitlera, wprowadzając NSDAP na scenę narodową. Potentaci biznesowi zapewnili nazistom hojne fundusze na floty samochodów i głośników, aby nasycić ulice niemieckich miast i wsi oraz wspierali finansowo organizacje powiązane z partią nazistowską, grupy młodzieżowe i siły paramilitarne. W trakcie kampanii w lipcu 1932 roku, Hitler miał wystarczające fundusze, żeby odwiedzić 50 miast w zaledwie 2 tygodnie.

W trakcie tej kampanii naziści otrzymali 37.3\% głosów - był to najwyższy uzyskany przez nich do tej pory wynik w demokratycznych wyborach narodowych. Nigdy nie mieli większości społeczeństwa po swojej stronie. Jeśli mieli jakąkolwiek wiarygodną bazę poparcia, to zazwyczaj wśród zamożniejszej części społeczeństwa. Ponadto, elementy drobnej burżuazji i wielu lumpenproletariuszy służyło jako partyjne zbiry, zorganizowane w oddziały szturmowe SA. Lecz zdecydowana większość zorganizowanej klasy robotniczej do samego końca wspierała komunistów lub socjaldemokratów.

W grudniowych wyborach prezydenckich 1932 roku startowało trzech kandydatów: konserwatywny urzędujący feldmarszałek von Hindenburg, nazistowski kandydat Adolf Hitler i kandydat partii komunistycznej Ernst Thälmann. W swojej kampanii Thälmann zapewniał, że głos na Hindenburga był równoznaczny z głosem na Hitlera, a Hitler zaprowadziłby Niemcy na wojnę. Kapitalistyczna prasa, włącznie z socjaldemokratami, potępiła ten pogląd jako "inspirowany Moskwą". Hindenburg został ponownie wybrany, w czasie gdy naziści stracili ok. 2 milionów głosów w wyborach do Reichstagu, w porównaniu do ich rekordowego wyniku, który wynosił 13.7 miliona głosów.

Zgodnie z przewidywaniami, socjaldemokratyczni przywódcy odrzucili propozycję partii komunistycznej, by w ostatniej chwili zawiązać koalicję przeciwko nazistom. Jak w wielu przypadkach, wcześniej i później, socjaldemokraci prędzej sprzymierzyliby się z reakcyjną prawicą niż znaleźliby wspólny język z Czerwonymi.\footnote{Wcześniej, w 1924 roku, socjaldemokratyczni przywódcy w Ministerstwie Spraw Wewnętrznych wykorzystali faszystowskie bojówki Reichswehry i Freikorpsu, by zaatakować lewicowych demonstrantów. Uwięzli ponad tysiąc robotników i zdławili gazety partii komunistycznej: Richard Plant, \textit{The Pink Triangle}, 47.} W tym samym czasie, szereg partii prawicowych poparł nazistów i już w styczniu 1933 roku, zaledwie kilka tygodni po wyborach, Hindenburg zaoferował Hitlerowi stanowisko kanclerza.

Po przejęciu władzy państwowej, Hitler i jego poplecznicy realizowali program polityczno-ekonomiczny podobny do programu faszystów włoskich. Rozbili zorganizowane związki zawodowe i pozbyli się wszelkich wyborów, partii opozycyjnych i niezależnych publikacji. Uwięzili, torturowali lub zamordowali setki tysięcy przeciwników politycznych. W Niemczech, podobnie jak we Włoszech, komuniści spotkali się z najcięższymi represjami politycznymi.

I oto dwa narody, Włosi i Niemcy, o odmiennej historii, kulturze, języku i rzekomo odmiennym usposobieniu, zastosowały te same represyjne metody z powodu uderzającego podobieństwa władzy ekonomicznej i konfliktu klasowego panującego w ich państwach. W tak różnorodnych państwach jak Litwa, Chorwacja, Rumunia, Węgry czy Hiszpania, pojawił się podobny faszystowski schemta, który starał się uchronić  wielki kapitał przed narzuceniem mu demokracji.\footnote{Nie zmienia to faktu, że kulturowe różnice mogą prowadzić do istotnych rozbieżności. Przykładem może być choćby przerażająca rola antysemityzmu w nazistowskich Niemczech w porównaniu z faszystowskimi Włochami.}

\sekcja{Kogo wspierali faszyści?}

Jest mnóstwo literatury poświęconej temu, kto wspierał nazistów, ale relatywnie mało temu, kogo wspierali naziści po przejęciu władzy. Jest to w zgodzie z tendencją konwencjonalnych prac, by unikać całkowicie tematu kapitalizmu, gdy cokolwiek negatywnego może być o nim powiedziane. Czyje interesy wspierali Mussolini i Hitler?

Zarówno we Włoszech w latach 20. jak i w Niemczech w latach 30., stare przemysłowe represje - o którcyh myślano, że już na zawsze odeszły w niepamięć - pojawiły się ponownie wraz z drstycznym pogorszeniem się warunków pracy. W imieniu ocalenia świata przed \textit{czerwoną zarazą} zdelegalizowano związki zawodowe i strajki. Własność związków i kooperatyw zostałą skonfiskowana i przekazana bogatym właścicielom prywatnym. Płaca minimalna, wynagrodzenie za nadgodizny i regulacje bezpieczeństwa zostały zniesione. ... Zwolnienia i aresztowania czekały na pracowników, którzy narzekali na niebezieczne lub nieludzkie warunki pracy. Już i tak skromne płace zostały drastycznie obniżone: w Niemczech o 25-40\%, a we Włoszech o 50\%. We Włoszech przywrócono pracę dzieci.

Oczywiście rzucono parę okruchów masom. Były darmowe koncerty i wydarzenia sportowe, mizerne programy społeczne, zasiłek dla bezrobotnych finansoweany głównie z wkładu ludzi pracujących i pokazowe publiczne prace mające na celu wzbudzić dumę narodową.

Zarówno Mussolini jak i Hitler odwdzięczyli się swoim biznesowym patronom prywatyzując wiele rentownych państwowych hut stali, elektrowni i spółek zajmujących się produkcją parowców. Oba reżimy sięgały chętnie do skarbu państwa, aby ratować lub dotować przemysł ciężki. Rolnictwo agrobiznesowe było rozwijane i masywnie dotowane. Oba państwa gwarantowały zwrot kapitału zainwestowanego przez wielkie korporacje, samemu ponosząc większość ryzyka i strat spowodowanych niekorzystnymi inwestycjami. Jak to nierzadko bywa z reakcyjnymi reżimami, publiczny kapitał został ograbiony przez prywatny.

W tym samym czasie podniesiono podatki dla pracowników, jednocześnie obniżając lub eliminując je dla bogatych i wielkiego biznesu. Podatek od spadku dla bogatych został znacznie zmniejszony lub całkowicie zlikwidowany.

Jakie były tego skutki? W latach 30. we Włoszech gospodarka była dotknięta recesją, oszałamiającym długiem publicznym i wszechobecną korupcją. Ale zyski w przemyśle rosły, a fabryki zbrojeniowe intensywnie produkowały broń w ramach przygotowań do nadchodzącej wojny. W Niemczech bezrobocie spadło o połowę, na skutek pokaźnej ekspansji w zatrudnieniu w przemyśle zbrojeniowym, ale ogólne ubóstwo rosło na skutek drastycznych obniżek płac. W latach 1935-43 zyski w przemyśle pokaźnie wzrosły w czasie, gdy dochód korporacyjnych przywódców wzrosł o 46\%. W czasie kryzysu lat 30. w USA, Wielkiej Brytanii i Skandynawii zarobki grup o wyższych dochodach uległy niewielkiemu spadkowi w udziale dochodu państwowego; lecz w Niemczech 5\% najbogatszych odnotowało 15-procentowy wzrost.\footnote{Simon Kuznets, "Qualitative Aspects of the Economic Growth of Nations", \textit{Economic Development and Cultural Change}, 5, no. 1, 1956, 5-94.}

Pomimo tego rekordu, większość autorów ignoruje bliską współpracę faszyzmu z wielkim biznesem. Niektórzy twierdzą nawet, że biznes nie był beneficjentem, lecz ofiarą faszyzmu! Angelo Codevilla, konserwatywny autor \textit{Hoover Institution} beztrosko oznajmił: "jeżeli faszyzm oznacza cokolwiek, to oznacza rządową własność i kontrolę nad biznesem" (\textit{Commentary, 8/94}). Zatem faszyzm jest nam prezentowany jako zmutowana forma socjalizmu. W rzeczywistości, jeżeli faszyzm oznacz cokolwiek, to oznacza pełne rządowe wsparcie dla biznesu i bezwzględne represje wobec ruchu roboticzego.\footnote{Były lewicowiec i odrodzony konserwatysta, Eugene Genovese (\textit{New Republic, 4/1/95}) z zapałem doszedł do wnosku, że "uznanie faszyzmu za twór wielkiego kapitału" jest "bezsensowną interpretacją". Genovese przyklaskiwał Ericowi Hobsbawmowi, który twierdził, że klasa kapitalistów nie była główną siłą stojącą za faszymem w Hiszpanii. W odpowiedzi Vicente Navarro (\textit{Monthly Review 1/96 oraz 4/96}) zauważył, że "główne interesy ekonomiczne Hiszpanii", wspierane przez co najmniej jednego teksańskiego milionera naftowego i inne elementy międzynarodowego kapitału faktycznie finansowały faszystowską inwazję i zamach stanu dokonane przez Franco przeciwko Republice Hiszpańskiej. Kluczowym źródłem, jak pisze Navarro, było imperium finansowe Joana Marcha, założyciela partii liberalnej i właściciela liberalnej gazety. Uważany za modernizatora i alternatywę dla oligarchicznego, opartego o ziemię, reakcyjnego sektora kapitału, March znalazł wspólny cel z tymi samymi oligarchami, gdy zrozumiał, że partie robotnicze zyskiwały poparcie i jego własne interesy ucierpiały na skutek reformistycznej politykii Republiki.} Czy faszyzm jest zatem po prostu dyktatorczą siłą w służbie kapitału? Nie jest to może jego całościowa charakterystyka, ale niewątpliwie stanowi to istotną część racji bytu faszyzmu, funkcją, o której sam Hitler wspominał kiedy mówił o ratowaniu industrialistów i bankowców przed bolszewizmem. Jest to zagadnienie, które zasługuje na zdecydowanie więcej uwagi niż mu poświęcono.

Mimo że faszyści rzeczywiście mogli wierzyć w to, że ratują plutokratów przed czerwonymi, w rzeczywistości rewolucyjna lewica nigdy nie była wystarczająco silna, by przejąć władzę państwową, zarówno w Niemczech, jak i we Włoszech. Siły ludowe były jednak wystarczająco silne, by obniżyć stopy zysku i interweniować w procesie akumulacji kapitału. To frustrowało kapitalistów próbujących rozwiązać wewnętrzne sprzeczności obciążając coraz większymi kosztami lud pracujący. Rewolucyjny czy nie, ten demokratyczny ruch robotniczy stanowił poważną przeszkodę dla interesów finansowych.

Oprócz służenia kapitalistom, faszystowscy przywódcy służyli samym sobie, czerpiąc zyski przy każdej nadarzającej się okazji. Ich osobista chciwość i ich lojalność klasowa były dwoma stronami tej samej monety. Mussolini i jego poplecznicy żyli wystawnie balująć z wyższymi kręgami bogaczy i arystokracji. Nazistowscy urzędnicy i dowódcy SS gromadzili osobiste fortuny rabując podbite tereny i okradając więźniów obozów koncentracyjnych i inne polityczne ofiary. Zarabiano ogromne sumy na potajemnych firmach z dobrymi znajomościami oraz z wynajmowania przemysłowym firmom takim jak \textit{I.G. Farben} albo \textit{Krupp} więźniów obozów koncentracyjnych do niewolniczej pracy.

Hitler jest zazwyczaj przedstawiany jako ideologiczny fanatyk, niezainteresowany przyziemnymi dobrami materialnymi. W rzeczywistości, zgromadził ogromną fortunę, w znacznej mierze zdobytą w wątpliwy sposób. Skonfiskował wiele dzieł sztuki, które wcześniej były własnością publiczną. Ukradł olbrzymie sumy pieniędzy z budżetu NSDAP. Stworzył koncept "prawa do osobowości", który pozwolił mu pobierać niewielką opłatę za każdy znaczek pocztowyz jego wizerunkiem, co przyniosło mu setki milionów marek.\footnote{Już przedtem istniał znaczek z wizerunkiem von Hindenburga mający na celu uhonorować jego prezydenturę. Stary Hindenburg, który nie darzył Hitlera specjalnym uwielbieniem, rzekł sarkastycznie, że uczyni Hitlera swoim ministrem poczty, bo "wtedy będzie mógł lizać mi tyłek".}

Największym źródłem bogactw Hitlera był sekretny fundusz łapówkowy, który najwięksi niemieccy industrialiści regularnie dotowali. Hitler "wiedział, że tak długo, jak niemiecki przemysł zarabia, jego prywatne źródła zarobku się nie wyczerpią. Zadbał zatem, żęby niemiecki przemysł nigdy nie miał się lepiej niż pod jego rządami - poprzez uruchamianie gigantycznych projektów zbrojeniowych",\footnote{Wulf Schwarzwaeller, \textit{The Unknown Hitler, 197}.} czyli tego co dzisiaj nazwalibyśmy wielkimi kontraktami obronnymi.

Hitler, daleki od ascetyzmu, prowadził życie pełne wygód i zachcianek. Przez cały czas sprawowania urzędu otrzymywał specjalne orzeczenia z niemieckiego urzędu skarbowego, które umożliwiały mu unikać płacenia podatku od przychodu i od nieruchomości. Posiadał całą flotę limuzyn, prywatne apartamenty, domy letniskowe, liczną służbę i majestatyczną posiadłość w Alpach. Najszczęśliwsze chwile spędzał goszcząc członków europejskich rodzin królewskich, w tym księcia i księżną Windsoru, którzy należeli do jego entuzjastycznych wielbicieli.

\sekcja{Gratulacje dla Adolfa i Benito}
Włoski faszyzm i niemiecki nazizm znalazły wielu wielbicieli w społeczności biznesowej Stanów Zjednoczonych i w korporacyjnej prasie. Bankowcy, wydawcy i industrialiści, wliczając w to osoby pokroju Henry'ego Forda, przybywali do Rzymu i Berlina, by składać hołd, odbierać odznaczenia i zawierać korzystne umowy. Wielu z nich zrobiło co mogło, by wspomóc nazistowskie zbrojenia wojenne, dzieląc się z nazistami militarno-przemysłowymi sekretami i dokonując z nimi potajemnych umów, nawet po tym, gdy Stany Zjednoczone dołączyly do wojny.\footnote{Charles Higham, \textit{Trading with the Enemy} (Nowy Jork: Dell, 1983).} W latach 20. i wczesnych latach 30. wiele znaczących publikacji, takich jak \textit{Fortune}, \textit{The Wall Street Journal}, \textit{Saturday Evening Post}, \textit{New York Times}, \textit{Chicago Tribune} czy \textit{Christian Science Monitor} chwaliło Mussoliniego jako człowieka, który ocalił Włochy od anarchii i radykalizmu. Snuli pełne zachwytu fantazje o odrodzonych Włoszech, gdzie nędza i wyzysk nagle zniknęły, gdzie czerwoni zostali pokonani, zapanowała harmonia, a Czarne Koszule chroniły "nową demokrację".

Włoskojęzyczna prasa w Stanach Zjednoczonych chętnie dołączyłą do tego chóru. Dwie najbardziej wpływowe gazety, \textit{L'Italia} wydawana w San Francisco, finansowana głównie przez Bank of America należący do A.P. Gianniniego oraz \textit{Il Progresso} wydawana w Nowym Jorku, której właścicielem był multimilioner Generoso Pope, patrzyły przychylnie na faszystowski reżim i zapewniały, że Stany Zjednoczone mogłyby skorzystać z podobnego porządku społecznego. Niektórzy odszczepieńcy odmówili przyłączenia się do chóru "Uwielbiamy Benito". \textit{The Nation} przypominało swoim czytelnikom, że Mussolini nie \textit{ratował} demokracji, tylko ją \textit{niszczył}, lecz ich krytyczne sentymenty nie zyskały popularności w korporacyjnych mediach Stanów Zjednoczonych.

Podobnie sprawa miała się z Hitlerem. Prasa rzadko ukazywała nazistowską dyktaturę \textit{Führera} w złym świetle. Istniała silna grupa twierdząca, że "powinniśmy dać Adolfowi szansę", częściowo wspierana przez nazistowskie fundusze. Na przykład za bardziej pozytywne artykuły w gazetach Hearsta, naziści płacili prawie dziesięciokrotnie więcej niż wynosiła standardowa opłata. W zamian William Randolph Hearst polecił swoim korespondentom w Niemczech, aby składali przychylne raporty na temat rządóœ Hitlera. Ci, którzy odmówili byli przenoszeni lub zwalniani. W gazetach Hearsta gościły nawet od czasu do czasu gościnne felietony wybitnych nazistowskich przywódców, takich jak Alfred Rosenberg i Hermann Göring.

W drugiej połowie lat 30. Włochy i Niemcy w sojuszu z Japonią, kolejnym dopiero co rozwijającym swój przemysł państwie, agresywnie dążyły do zdobycia udziału w światowych rynkach i łupach kolonialnych. Ten ekspansjonizm wpędzał je w coraz większe konflikty z bardziej rozwiniętymi zachodnimi państwami kapitalistycznymi, takimi jak Wielka Brytania, Francja i Stany Zjednoczone. W miarę narastania wojennego napięcia, opinia amerykańskiej prasy na temat państw Osi przybrała zdecydowanie bardziej krytyczny ton.

\sekcja{Racjonalne wykorzystanie irracjonalnej ideologii}
Niektórzy autorzy szczególnie uwydatniają "irracjonalne" cechy faszyzmu. Czyniąc to, pomijają racjonalne funkcje polityczno-ekonomiczne, jakie spełniał faszyzm. W dużej mierze polityka to racjonalna manipulacja irracjonalnymi symbolami. Z pewnością dotyczy to ideologii faszystowskiej, której emocjonalny wydźwięk pełni funkcję kontroli klasowej.

Zaczęło się od kultu przywódcy, we Włoszech: \textit{il Duce}, a w Niemczech: \textit{der Führerprinzip}. Ramię w ramie z kultem przywódcy szło ubóstwianie państwa. Jak pisał Mussolini \textit{"koncepcja faszystowska jest antyindywidualistyczną, staje wszakże na stanowisku jednostki, o ile ta utożsamia się z państwem"}. Faszyzm wyznaje autorytarną zasadę wszechogarniającego państwa i najwyższego przywódcy. Wychwala surowsze ludzkie impulsy podboju i dominacji, jednocześnie odrzucając egalitaryzm, demokrację, kolektywizm i pacyfizm jako doktryny słabości i dekadencji.

Poświęcenie pokojowi - pisał - jest "wrogie faszyzmowi". Ciągły pokój - jak twierdził w 1934 r. - jest "przygnębiającą" doktryną. Jedynie w "okrutnej walce" i "podboju" mężczyźni i narody osiągają swoje najwyższe spełnienie. "Mimo, że słowa to piękna rzecz" - zapewniał - "karabiny, pistolety maszynowe, samoloty i armaty są wciąż piękniejsze". Innym razem pisał: "Jedynie wojna wznosi do maksimum napięcia wszystkie siły ludzkie i znaczy znamieniem szlachectwa narody, które mają odwagę stawić jej czoło". Jak na ironię, większość włoskich poborowych nie miała ochoty na wojny Mussoliniego i miała zwyczaj wycofywać się z walki, gdy tylko odkrywała, że druga strona używa ostrej amunicji.

Faszystowska doktryna uwypukla monistyczne wartośći: \textit{Ein Volk}, \textit{ein Reich}, \textit{ein Führer} (jeden lud, jeden rząd, jeden przywódca). Ludzie nie mają się już przejmować podziałami klasowymi, lecz postrzegać siebie jako część harmonijnej całości, w której biedni i bogaci stanowią jedność - pogląd, który wspiera ekonomiczny statuts quo, maskując trwający system wyzysku klasowego. Stoi to w kontrze do programu lewicowego, który opowiada się za stawianiem żądań przez masy i uwydatnianiem niesprawiedliwości społecznej i walki klasowej.

Ten monizm jest wzmocniony atawistycznymi odwołaniami do mitycznych korzeni narodu. Dla Mussoliniego tą wielkością był starożyny Rzym, dla Hitlera - antyczny Volk. Sztuka napisana przez pro-nazistowskiego Hansa Jorsta, zatytułowana \textit{Schlageter} i masowo wystawiana w Niemczech po przejęciu władzy przez nazistów (Hitler był obecny na premierze w Berlinie), przeciwstawia mistycyzm Volku polityce klasowej. Pełen entuzjazmu August rozmawia ze swoim ojcem, Schneiderem:

\dialog{August}{Nie uwierzysz, Tato, ale \ldots młodzi ludzie zwracają już większej uwagi na te dawne hasła \ldots walka klasowa wygasa.}
\dialog{Schneider}{No więc, co macie zamiast tego?}
\dialog{August}{Społeczność Volk}
\dialog{Schneider}{I to jest hasło?}
\dialog{August}{Nie, to doświadczenie!}
\dialog{Schneider}{Mój Boże, nasza walka klasowa, nasze strajki, one nie były doświadczeniem, prawda? Socjalizm, Międzynarodówka, może były tylko fantazjami?}
\dialog{August}{Były konieczne, ale \ldots są historycznymi doświadczeniami.}
\dialog{Schneider}{A zatem wasza wspólnota Volk jest przyszłością. Powiedz mi, jak ją sobie właściwie wyobrażasz? Biedni, bogaci, zdrowi, klasa wyższa i niższa, wszystko to u was znika, co? \ldots}
\dialog{August}{Słuchaj tato, klasa wyższa i niższa, biedni, bogaci - to zawsze istnieje. Decydujące jest tylko waga, jaką przywiązujemy do tego pytania. Dla nas życie nie jest podzielone na godziny pracy i opatrzone wykresami cen. Wierzymy raczej w ludzką egzystencję jako całość. Nikt z nas nie uważa zarabiania pieniędzy za rzecz najważniejszą; chcemy służyć. Jenostka jest częścią obiegu swojego ludu.\footnote{George Mosse (ed.), \textit{Nazi Culture} (Nowy Jork: Grosset \& Dunlap, 1966), 116-118.}}

Komentarze syna ujawniają: "walka klasowa wygasa".Obawy taty odnośnie nadużyć władzy klasowej i niesprawiedliwości klasowej są pochopnie odrzucane jako stan umysłu bez odzwierciedlenia w rzeczywistości. Fałszywie utożsamia się je nawet z troską o pieniądze. ("Nikt z nas nie uważa zarabiania pieniędzy za rzecz najważniejszą"). Można przypuszczać, że sprawy bogactwa należy pozostawić tym, którzy je posiadają. Mamy coś lepszego, twierdzi August: totalistyczne, monistyczne doświadczenie jako naród, my wszyscy, bogaci i biedni, pracujące razem dla jakiejś większej chwały. Wygodnie pomija się, że "chwalebne ofiary" są ponoszone przez biednych dla dobra bogatych.

Pogląd wyrażony w tej sztuce i innych nazistowskich utworach propagandowych nie ujawnia obojętności wobec klas; wręcz przeciwnie ukazuje głęboką świadomość klasowych interesów, dobrze wykonaną próbę zamaskowania i stłumienia silnej świadomości klasowej, jaka istniała wśród robotników w Niemczech. W przebiegłym zaprzeczeniu często znajdujemy ukryte przyznanie się do winy.

\sekcja{Patriarchat i pseudo-rewolucja}
\sekcja{Przyjaźni dla faszyzmu}

\chapter{Chwalmy teraz rewolucję}
\sekcja{Koszty kontrrewolucji}
\sekcja{Czyja przemoc?}
\sekcja{Wolny rynek dla nielicznych}
\sekcja{Wolność rewolucji}
\sekcja{Jaka miara bólu?}

\chapter{Lewicowy antykomunizm}
\sekcja{Ugięcie się przed ortodoksją}
\sekcja{Czysty socjalizm, a socjalizm w oblężeniu}
\sekcja{Decentralizacja, a przetrwanie}

\chapter{Komunizm w Krainie Czarów}
\chapter{Palce Stalina}

\end{document}
